\section{递归神经张量网络与情感分析}

\subsection{情感分析的挑战}

情感分析看似简单——检测正面词（great, wonderful）和负面词（poor, bad）即可。但真实语言中的情感表达远比这复杂：

\begin{figure}[H]
\centering
\includegraphics[width=0.95\textwidth]{slides-images/slide-051.jpg}
\caption{情感分析的难点示例：``With this cast and subject matter, the movie should have been funnier and more entertaining''——虽然包含正面词，但整体是负面评价\protect\footnotemark}
\end{figure}
\footnotetext{Slides 第52页。}

\begin{warningbox}{关键词匹配的失败案例}
考虑影评``With this cast and subject matter, the movie should have been funnier and more entertaining.'' 关键词匹配会检测到 ``funnier''（正面）和 ``entertaining''（正面），将其判为正面评价。但由于 ``should have been'' 的组合语义，这实际上是一条负面评价——影评人认为这部电影\textbf{不够}有趣和娱乐性。这种组合语义是简单方法无法捕获的。
\end{warningbox}

\subsection{Stanford Sentiment Treebank}

为了研究组合语义在情感分析中的作用，Manning 团队构建了 \textbf{Stanford Sentiment Treebank（SST）}：

\begin{figure}[H]
\centering
\includegraphics[width=0.95\textwidth]{slides-images/slide-050.jpg}
\caption{Stanford Sentiment Treebank：对近 12,000 个句子的每个短语节点都标注了情感标签\protect\footnotemark}
\end{figure}
\footnotetext{Slides 第51页。}

\begin{importantbox}{SST 数据集的创新}
\begin{itemize}
    \item 包含 \textbf{11,855} 个句子，共 \textbf{215,154} 个带情感标注的短语
    \item 不仅标注整句情感，还标注了\textbf{每一个中间短语}的情感
    \item 使用 5 级情感标注：非常负面、略微负面、中性、略微正面、非常正面
    \item 仅仅使用 phrase-level 标注就能提升 Naive Bayes 分类器 4\% 的准确率（79\% $\rightarrow$ 83\%）
\end{itemize}
\end{importantbox}

\subsection{递归神经张量网络（RNTN）}

为了克服简单 TreeRNN 的表达能力限制，Manning 团队提出了\textbf{递归神经张量网络（Recursive Neural Tensor Network, RNTN）}。

\begin{figure}[H]
\centering
\includegraphics[width=0.95\textwidth]{slides-images/slide-055.jpg}
\caption{RNTN 的 Neural Tensor Layer：通过张量实现词对之间的乘法交互\protect\footnotemark}
\end{figure}
\footnotetext{Slides 第56页。}

RNTN 的核心创新是 \textbf{Neural Tensor Layer}：

$$\mathbf{p} = f\left(\begin{bmatrix} \mathbf{c}_1 \\ \mathbf{c}_2 \end{bmatrix}^T V^{[1:d]} \begin{bmatrix} \mathbf{c}_1 \\ \mathbf{c}_2 \end{bmatrix} + W \begin{bmatrix} \mathbf{c}_1 \\ \mathbf{c}_2 \end{bmatrix} + \mathbf{b}\right)$$

其中：
\begin{itemize}
    \item $V^{[1:d]} \in \mathbb{R}^{2d \times 2d \times d}$ 是三维张量
    \item $\mathbf{c}_1^T V^{[k]} \mathbf{c}_2$ 计算了两个子节点之间的\textbf{乘法交互}
    \item $W$ 项保留了标准的线性组合
    \item 张量项使得模型能根据两个子节点的具体内容，产生不同的组合方式
\end{itemize}

\begin{knowledgebox}{为什么需要乘法交互}
在语言中，不同类型的词组合方式差异很大：
\begin{itemize}
    \item ``red ball''（形容词+名词）：red 为 ball 添加属性
    \item ``kick the ball''（动词+宾语）：ball 是动作的对象，角色完全不同
    \item ``not good''（否定+正面词）：not 翻转了 good 的极性
\end{itemize}
简单的线性组合 $W[\mathbf{c}_1; \mathbf{c}_2]$ 无法区分这些不同的交互模式。张量层通过 $\mathbf{c}_1^T V \mathbf{c}_2$ 实现了\textbf{依赖于输入内容}的组合方式——当 $\mathbf{c}_1$ 和 $\mathbf{c}_2$ 不同时，张量项会产生不同的交互效果。
\end{knowledgebox}

\subsection{RNTN 在情感分析上的表现}

\begin{figure}[H]
\centering
\includegraphics[width=0.95\textwidth]{slides-images/slide-053.jpg}
\caption{RNTN 准确率对比：红色曲线（RNTN）在各长度的短语上均优于 Naive Bayes 和早期 TreeRNN\protect\footnotemark}
\end{figure}
\footnotetext{Slides 第54页。}

RNTN 相比 bigram Naive Bayes 模型提升了约 2 个百分点。虽然绝对提升不大，但更有价值的是其对\textbf{组合语义}的建模能力。

\subsection{否定建模：RNTN 的亮点}

RNTN 最令人印象深刻的能力是对\textbf{否定（negation）}的处理：

\begin{figure}[H]
\centering
\includegraphics[width=0.95\textwidth]{slides-images/slide-057.jpg}
\caption{RNTN 处理否定：``It's just incredibly dull'' vs ``It's definitely not dull''——模型正确捕获了否定对情感极性的翻转效果\protect\footnotemark}
\end{figure}
\footnotetext{Slides 第58页。}

\begin{importantbox}{否定建模的关键测试}
考虑两种情况：
\begin{enumerate}
    \item \textbf{否定正面句}（如``not good''）：相对容易，因为 ``not'' 本身在统计上与负面语境相关，它的出现会自然降低正面程度
    \item \textbf{否定负面句}（如``not dull''）：这是\textbf{真正的挑战}。``not'' 是负面词，``dull'' 也是负面词，简单的加法/平均模型会得出``更加负面''的错误结论。正确答案应该是情感被\textbf{翻转为正面}
\end{enumerate}
RNTN 是当时唯一能正确处理``否定负面词使情感变正面''这一现象的模型。这种能力来自张量层的乘法交互——它能学会``not + 负面词 = 正面''这样的非线性组合规则。
\end{importantbox}

\begin{figure}[H]
\centering
\includegraphics[width=0.95\textwidth]{slides-images/slide-058.jpg}
\caption{否定建模的定量分析：RNTN（绿色）在否定负面句上正确地将情感推向正面，而其他模型无法做到\protect\footnotemark}
\end{figure}
\footnotetext{Slides 第59页。}

\subsection{可解释的组合语义}

RNTN 的另一个优点是其\textbf{可解释性}。由于模型在句法树的每个节点都产生情感预测，我们可以观察情感如何在组合过程中变化：

\begin{figure}[H]
\centering
\includegraphics[width=0.95\textwidth]{slides-images/slide-056.jpg}
\caption{RNTN 在整棵句法树上的情感预测：可以清楚看到``slow and repetitive parts''（负面）和``enough spice to keep it interesting''（正面）如何组合为整体正面评价\protect\footnotemark}
\end{figure}
\footnotetext{Slides 第57页。}

这种逐节点的情感可视化在语言学上非常有意义——它展示了模型如何理解语言的组合语义。

\subsection{本章小结}

RNTN 通过张量层实现了词对之间的乘法交互，使得不同类型的语法组合可以产生不同的效果。在 Stanford Sentiment Treebank 上，RNTN 展现了优秀的组合语义建模能力，特别是在处理否定等复杂语言现象时。其逐节点的情感预测提供了很好的可解释性。

%% ========== 树递归网络的反思 ==========

